\chapter{Architektur und nachvollziehbare Beiträge}
\label{api:development}\index{Entwicklung}\index{TDD}\index{Assembly}
\begingroup\raggedright
\section{Vom Dokument zum prüfbaren Ergebnis}
Model liest/konvertiert dokumentierte Inputs und Regionen, Preflight prüft
Lagerung/Gauge/Elektroden. Mesh kennt Topologie und Sätze, Geometrie/Quadratur
kennen Tetraeder und Flächen. Materialien enthalten validierte SI-Tensoren.
Felder/DofMap legen Reihenfolge fest; Elemente/Physik erzeugen lokale
Matrix-/Lastbeiträge. Assembly bildet COO und CSR/CSC. Boundary reduziert
Dirichlet/Elektroden; Solver wählen Faktorisierung/Zeitverfahren und prüfen
feldweise Residuen/Bilanzen. Ergebnisobjekte sind owned Daten, IO/Plot
Adapter mit expliziten Form-/Maskenverträgen. Study/UQ orchestrieren diese
Pipeline, ändern keine wissenschaftlichen Operatoren. Surrogate verwenden
prüfbare Tabellen und eigene Trainings-/Testprovenienz.

Interne Erweiterungsgrenzen werden im Inventar klar benannt. Neue Elemente
brauchen Formfunktionen, Geometrie/Jacobi, Quadratur, lokale Operatoren,
Last-/Randintegrale, Postprocessing und IO-Formverträge; eine einzige neue
Steifigkeitsroutine genügt nicht. Neue Physik braucht Felddefinitionen,
Vorzeichen/SI, Nullraum-/BC-Vertrag, unabhängige Einzeltests und erst dann
Kopplung. Backends sehen bereits reduzierte Systeme mit Feldpartition;
physische Reduktion bleibt beim Solver. IO-/Visualisierungsadapter dürfen
Rohfelder nicht still glätten oder komplexe Komponenten verlieren.

\section{Kleine lokale Erweiterung mit unabhängigem Test}
Das folgende Lehrbeispiel implementiert einen lokalen linearen Wärmequellen-
Callback, keinen neuen produktiven Solver. Es nutzt ausdrücklich die
Forschungs-Lastgrenze und begründet ihren Vertrag: Punkt m, Quelle W/m$^3$,
Gradient W/m$^4$. Für das Einheitstetraeder mit $Q=2+3x$ sind
$\int N_iN_j\,dV=V/20$ ($i\ne j$) und $V/10$ ($i=j$).
Daraus folgt Lastvektor $[13/120,2/15,13/120,13/120]$ W und Summe $11/24$ W.
Die absolute Grenze $2\cdot10^{-15}$ W deckt Rundung der wenigen
Quadraturoperationen ab, keine numerisch gelernte Referenz.
\lstinputlisting{examples/developer_extension.py}
\begin{lstlisting}
python -I examples/developer_extension.py
PYTHONPATH=src python -m pytest tests/unit/test_developer_handbook_example.py
\end{lstlisting}
Der Archivaufruf braucht ein separat installiertes Kandidatenpaket; der
Testaufruf gehört zum Checkout. Der Test wurde zuerst am fehlenden Callback
rot ausgeführt, dann bestanden Referenzintegration und ungültiger Gradient.
Der Callback ist nur Lehre, keine neue stabile Produktions-API.

\section{TDD, Teststruktur und wissenschaftliche Toleranzen}
Entwicklungsinstallation: CPython 3.11--3.13, virtuelle Umgebung, Paket mit
\texttt{dev}-Extra; optionale Visualisierung/PETSc gezielt und getrennt.
Unit-Tests prüfen lokal Formen, ungültige Daten, Tensoren und Semantik.
Patch-Tests prüfen reproduzierbare konstante/affine Felder; Regressionen
binden tatsächliche Workflows und veröffentlichte Verträge. Convergence
nutzt mehrere Netze/Zeitschritte und unabhängige Referenzen; Backendtests
laufen zusätzlich in realen nativen PETSc-Umgebungen. Benchmarks messen
Zeit/Genauigkeit mit Hardware und Problemgröße, ohne universelle Zusage.

Neue ausführbare Funktion zuerst mit kleinem deterministischem Red-Test,
Green mit minimal korrekter Implementierung, Refactor bei unveränderter
Erwartung. Rot muss fehlendes Verhalten zeigen, keinen kaputten Testaufbau.
Keine künstlichen roten Tests für reine Prosa. Numerische Referenzen aus
Analytik, unabhängiger Diskretisierung oder Energie/Bilanz/Reziprozität;
absolute und relative Grenzen mit SI, Konditionierung und Rechengröße
begründen. Kein blindes Neuschreiben von Goldwerten, Abschalten von Tests
oder nachträgliches Aufweiten zur Maskierung eines Fehlers.
\begin{lstlisting}
python -m pytest
python -m ruff check .
python -m ruff format --check .
python -m mypy src/fem_toolbox
python scripts/build_handbook.py --check-matrix
python scripts/build_handbook.py --self-test
git diff --check
\end{lstlisting}
PETSc-lokale Skips sind ausdrücklich zu nennen und kein bestandener nativer
Nachweis. API-Driftprüfung läuft ohne Import optionaler Bibliotheken;
Signaturänderung, fehlendes Modul/Reexport oder fehlende Referenz werden abgelehnt.
Die Extraktion verwendet die Python-AST-Schnittstelle\cite{python-ast};
das Dokumentensystem ist LaTeX\cite{latex-project}.
\texttt{api-inventory.json} und daraus erzeugtes Inventarkapitel müssen zum
Quellstand passen. Manuelle Fachverträge bleiben Pflicht, auch bei grüner AST-Prüfung.

\section{Issue, Milestone, Review und Release}
Vor Implementierung passendes Issue mit Ziel, Umfang, Nichtzielen,
Akzeptanzkriterien und Prüfplan. Issue-Branch
\texttt{codex/<nummer>-<beschreibung>} vom Milestone-Zweig; PR zurück in
\texttt{codex/milestone-m9}, nicht direkt main. PR nennt Problem, Verhalten,
Red/Green und tatsächliche Prüflücken. Nur vollständige Arbeitspakete schließen;
fehlende unabhängige Reviewer offen benennen, keine fremde Freigabe erfinden.
Anwendbare CI am aktuellen Commit und geklärte Befunde vor Integration.

M9/main erst bei vollständigem Pflichtumfang, aktueller Verifikation und
Dokumentation. Aktuelles main vorher in M9 aufnehmen und den resultierenden
Stand prüfen. Nach Merge CI am tatsächlichen main-Zielcommit, saubere
Wheel-/sdist-Installation, alle Dokumentationsartefakte/Hashes/Commitmetadaten;
erst dann annotiertes unveränderliches Tag und Release. \#66 zuletzt schließen.
Experimentelle \#63/\#65 bleiben außerhalb M9 offen; keine Validierungsbehauptung
in Version, PDF oder Release Notes. Branchschutz ist gemäß ADR-020 kein
zusätzliches Gate, die Integrationsverantwortung bleibt bestehen.
\endgroup
